\section{Fixed-direction trichotomy}\label{sec:fixed}

Set \(v=(a,b,c,d)\in\R^4\) and \(H_v(\e)=H_0+\e D(v)\).  Define
\[
\ell=a-b+c-d,\qquad Q=bd-ac,\qquad
R=ab(c-d)+cd(a-b).
\]
The ray numerator is
\begin{equation}\label{eq:raynum}
N_v(z;\e)=\e\ell(z^2+2z)+2\e^2Q(z+1)+\e^3R.
\end{equation}

\begin{theorem}[Complete fixed-ray trichotomy]\label{thm:trichotomy}
Exactly one of the following holds.
\begin{enumerate}[label=(\alph*)]
\item If \(\ell\ne0\), the first perturbation/time pair is
      \((p,q)=(1,3)\), and the coefficient of \(\e t^3\) in \(K_v\) is
      \(i\ell/6\ne0\).
      Every other surviving exponent lies in
      \(((1,3)+\N^2)\setminus\{(1,3)\}\).
\item If \(\ell=0\) and \(Q\ne0\), then \((p,q)=(2,4)\), and the coefficient
      of \(\e^2t^4\) in \(K_v\) is \(Q/12\ne0\).
      Every other surviving exponent lies in
      \(((2,4)+\N^2)\setminus\{(2,4)\}\).
\item If \(\ell=Q=0\), then \(K_v(\e,t)\equiv0\); there is no cubic
      opening branch.
\end{enumerate}
\end{theorem}
\begin{proof}
The exact numerator \cref{eq:raynum} first separates the three cases.  If
\(\ell=Q=0\), the syzygy forces \(R=0\), so the rational entry and hence the
entire amplitude vanish.  For the first two cases, the endpoint support theorem
identifies the earliest possible time depth.  Direct boundary evaluation gives
\([\e]M_3=\ell\) and \([\e^2]M_4=2Q\).  Multiplication by
\((-i)^3/3!=i/6\) and \((-i)^4/4!=1/24\) gives the stated amplitude
coefficients.  Non-zero coefficients prove sharpness.
\end{proof}

\noindent\textbf{Formal status.}
The literal coefficient support and all three mutually exclusive branches are
\evidence{KERNEL-CHECKED} over complex directions by
\path{M3A.linearRayData_of_linear_ne},
\path{M3A.quadraticRayData_of_linear_zero_quadratic_ne},
\path{M3A.darkRayData_of_linear_quadratic_zero} and
\path{M3A.fixedRay_trichotomy}.  The claim that exactly one branch is selected
is additionally traced to \path{M3A.classifyRay_eq_linear_iff},
\path{M3A.classifyRay_eq_quadratic_iff},
\path{M3A.classifyRay_eq_dark_iff} and \path{M3A.classifyRay_spec}.  The exact
coefficient values are checked in M5, while M8 identifies the array with the
convergent matrix-exponential amplitude.  The real statement here follows by
scalar restriction.

The ordered pair \((p,q)\) is essential.  The numerator valuation \(p\) alone
does not encode time depth; \(q\) alone does not encode perturbative order.  For
the protected direction \(v_*=(1,0,-1,0)\), \(\ell=0\), \(Q=1\), and
the first amplitude coefficient is exactly (1/12).

\begin{corollary}[Dark rays]\label{cor:darkrays}
The fixed dark directions are precisely
\[
(a-b,c-d)=(0,0)\quad\text{or}\quad(a-d,c-b)=(0,0).
\]
\end{corollary}

\begin{remark}[Family-specific completeness]
The absence of a cubic branch is exact for this four-parameter diagonal family.
It is not a universal theorem for all analytic Hamiltonian perturbations.  A
two-state off-diagonal perturbation proportional to \(\e^3\) supplies an
immediate cubic-opening control.
\end{remark}
